\section{Uncertainty budget of the lattice rates}\label{sec:11.9}

A dark-matter particle of \SI{30}{MeV} with a heavy mediator and $\sigma_p = 10^{-40}\,\mathrm{cm^2}$ makes $5.26\times10^3$ events per kilogram-year in silicon above a threshold of \SI{100}{meV}. The number assumes a harmonic crystal, nuclei that scatter independently and the phonons of the PBE force constants of \cref{sec:5.2}. The free nuclear recoil of the standard rate gives $4.91\times10^3$ for the same threshold, and above \SI{400}{meV} it gives a sixteenth of the crystal rate. That difference is the error of the simpler calculation: the crystal response is the better one, and the free recoil is its limit at large $x_R$ (\cref{sec:11.4}). The uncertainty is how far the crystal number moves when the truncations of this chapter are lifted, and \cref{tab:11-nuisances} lists these steps in the classes and bases of \cref{tab:10-nuisances}. A next order computed with the force constants of \cref{sec:5.2} or with measured inputs is marked ``computed'', and its central value has already moved.

The exact $\omega$-integrated weight of \cref{eq:11.3-weight} moves the central value of each phonon order; what remains is the distribution of the coherent weight in $\omega$ inside each band, which \cref{eq:11.3-band} bounds at every point but does not fix. The residual is a tilt within each band, and it matters for a threshold inside a band at momentum transfers below a few $q_{\mathrm{BZ}}$. Above the impulse onset the zeroth and first moments of the recoil spectrum are exact on any energy surface, and the anharmonic template is the constrained tilt \cref{eq:8.8-tilt}, which keeps them. For soft-mode and light-atom crystals, and for solid helium near melting, neither a computed next order nor a small parameter exists, and the Markov--Krein band of \cref{eq:C.10-band} is the only available statement (\cref{lim:11.5c}): a range without a central value, whose edges a likelihood uses.

The phonons themselves are an input, not a truncation of the label, and their error plays the part of the spread between tabulations in \cref{sec:10.10}. Of the eight measured frequencies of silicon at $\mathrm{\Gamma}$, X and L, the PBE force constants of \cref{sec:5.2} miss by \SI{2.6}{\percent} on average and the zone-center optical mode by \SI{0.6}{\percent}. We take as template a uniform rescaling of all phonon frequencies by \SI{2.6}{\percent}. It moves the Poisson weights through $2W_\kappa \propto \langle\omega^{-1}\rangle$ and the impulse width through $\langle\omega\rangle^{1/2}$ together, and since the deviations of single modes have both signs, it is an upper template. The measured elastic tensors that fix the acoustic plateau and edge of \cref{sec:11.1} are room-temperature values, and the change on cooling has not been sized.

\begin{table}[tbp]
  \caption{Uncertainty budget of the lattice rates. Classes and bases as in \cref{tab:10-nuisances}. Sizes are for silicon unless another crystal is named; \cref{we:11-band30} is the silicon rate at \SI{30}{MeV} above \SI{100}{meV}, and $d$ the site difference of \cref{we:11-gaas}. Items without a row: the optical-phonon shift of \SI{0.15}{meV} (\cref{tab:8-nuisances}), the quantum-crystal and germanium-spectrum models (\cref{sec:11.4,sec:11.5}), the absorption line widths and resonator prescription (\cref{sec:11.8}), and the nuclear spin inputs.}
  \label{tab:11-nuisances}
  \centering\footnotesize\setlength{\tabcolsep}{3pt}
  \begin{tabular}{@{}>{\raggedright\arraybackslash}p{0.25\linewidth}>{\raggedright\arraybackslash}p{0.12\linewidth}>{\raggedright\arraybackslash}p{0.38\linewidth}>{\raggedright\arraybackslash}p{0.09\linewidth}>{\raggedright\arraybackslash}p{0.10\linewidth}@{}}
    \toprule
    Truncation or input & Type & Size & Basis & Section \\
    \midrule
    Pairs of nuclei dropped (incoherent approximation) & normali\-zation per order & within \SI{10}{\percent} above \SI{6.6}{keV} ($n = 2, 3$); \cref{we:11-band30}: $+2.6\times10^{-4}$ & computed & \cref{sec:11.3} \\
    Coherent weight inside a band & tilt & up to $2.7\times$ incoherent at \SI{10}{keV} & bound & \cref{sec:11.3} \\
    Phonon frequencies (PBE) & tilt & \SI{2.6}{\percent} mean, uniform rescale (upper estimate); \cref{we:11-band30}: $+0.24\,\%$ ($+3.1\,\%$ above \SI{400}{meV}) & spread & \cref{sec:5.2} \\
    Anharmonic recoil width, skewness ($n_A = 2$) & tilt (con\-strained) & kinetic energy within $\pm0.2\,\%$; third cumulant $\sim\eta_{\mathrm{anh}}^2$ & model; est. & \cref{sec:11.5} \\
    Anharmonic few-phonon weights ($n_A = 2$) & tilt in $q$ & $n = 2$: $+\SI{50}{\percent}$, $n = 4$: $27\times$ at $q_{\mathrm{BZ}}$; harmonic to \SI{10}{\percent} above 5, 12, \SI{20}{keV} & model & \cref{sec:11.5} \\
    Strong anharmonicity (soft modes, light atoms, helium) & any & Markov--Krein band; Si, \SI{100}{keV}: 0.19 to 0.86 at $E_R$ & bound & \cref{sec:11.5} \\
    Thermal diffuse scattering ($n_T = 2$) & new support & $\propto T^2$; $10^{-9}$ of room temperature at \SI{10}{mK} & model & \cref{sec:11.3} \\
    Charge vertex above the zone & tilt in $q$ & unknown at 3--\SI{20}{keV}; nucleon coupling unaffected & named & \cref{sec:11.3} \\
    Free recoil for the crystal (evaluation error) & tilt & \SI{10}{\percent} from $x_R = 2.3$ (\SI{1}{GeV}); \cref{we:11-band30}: $-6.8\,\%$, $\times0.062$ above \SI{400}{meV} & computed & \cref{sec:11.4} \\
    Sound speeds & normali\-zation; threshold & computed against measured tensors & spread & \cref{sec:17.5} \\
    $Z^*/\varepsilon_\infty$, polar rates & normali\-zation & splitting route: rate \SI{2}{\percent}; semilocal $\varepsilon_\infty$: $-37\,\%$ & spread & \cref{sec:11.2} \\
    Charge of acoustic waves & new support & at most \SI{3}{\percent} near the edge, \SI{1}{\percent} at \SI{100}{keV} (GaAs); quadrupoles not computed & est. & \cref{sec:11.2} \\
    $Z^*/\varepsilon_\infty$, GaAs scalar absorption & normali\-zation & \SI{1}{\percent} gives \SI{9}{\percent}; PBE inputs: $/250$ & spread & \cref{sec:11.8} \\
    Core kinetic energy; mean masses (GaAs scalar) & normali\-zation & $\Delta d \approx -0.1\,d$, rate $-\SI{19}{\percent}$ (parametric); \SI{0.01}{u}: \SI{5}{\percent} of $d$ & est. & \cref{sec:11.8} \\
    Single-isotope density of states & normali\-zation & $+\SI{8.0}{\percent}$ at \SI{50}{meV} & computed & \cref{sec:11.6} \\
    Cold in-band absorption of silicon & normali\-zation & $\varGamma_{\mathrm{iso}}/\varGamma_{\mathrm{EM}} \ge 5.0\times10^3$; warm datum \SI{10}{\percent} & bound & \cref{sec:11.6} \\
    Electronic cubic harmonics & tilt in time & Si $A_{\mathrm{day}} = \SI{0.9}{\percent}$ without local fields & model & \cref{sec:11.7} \\
    \bottomrule
  \end{tabular}
\end{table}

A truncation that two rates share carries one parameter in both (\cref{sec:8.8}). The phonon frequencies of one set of force constants enter the one-phonon, multiphonon and impulse rates of that crystal together, at every mass and threshold, and a fit that gave each its own $\vartheta$ would let the data trade them. The measured elastic tensor of \cref{sec:11.1} is a different input, and an acoustic rate built from it shares no parameter with the multiphonon rate built from computed phonons. The ratio $Z^*/\varepsilon_\infty$ of a polar crystal is one input of its dark-photon scattering rate, its dark-photon absorption and its electron-scalar absorption, and one $\vartheta$ moves all three, by different factors. The halo stays fixed, as in \cref{eq:8.8-like}, and the detector uncertainties belong to \cref{sec:21.4}.

\begin{workedexample}{The uncertainty band of the silicon multiphonon rate at \SI{30}{MeV}}\label{we:11-band30}
Take $\mchi = \SI{30}{MeV}$, a heavy mediator, equal couplings to protons and neutrons with $\sigma_p = 10^{-40}\,\mathrm{cm^2}$, and the standard halo of \cref{sec:1.2} with $\rho_\chi = \SI{0.3}{GeV.cm^{-3}}$. The threshold $\omega_{\mathrm{th}} = \SI{100}{meV}$ and the cross-section $\sigma_p$ are illustrative choices; the threshold lies above the one-phonon band, so that every event counted contains at least two phonons. The largest momentum transfer is $2\mchi v_{\mathrm{max}} = \SI{159}{keV}$, where $E_R = \SI{0.48}{eV}$ and $x_R = 11.8$. The central value is \cref{eq:8.2-spectrum} with one row and the incoherent harmonic response of all orders at $T = 0$, from the PBE density of states of \cref{sec:5.2}, the Helm form factor, and the velocity integral done for an orientation-averaged crystal:
\begin{equation*}
  \mathcal{R}^{\mathfrak{T}} = \SI{5.26e3}{kg^{-1}.yr^{-1}} ,
\end{equation*}
with $\mathcal{R}^{\mathfrak{T}}$ the rate at the central values of the truncation labels (\cref{sec:10.10}), or 14 events per kilogram and day. Half of it comes from momentum transfers below \SI{80}{keV}, where $x_R = 3.0$, and \SI{80}{\percent} from between 54 and \SI{109}{keV}, $x_R$ from 1.4 to 5.6: the rate spans the impulse onset.

The computed coherent step multiplies each phonon order by $1 + \delta^{\mathrm{coh}}_n(q)$ from \cref{eq:11.3-weight}. Between 30 and \SI{159}{keV}, $|\delta^{\mathrm{coh}}_2|$ stays below 0.021 and $|\delta^{\mathrm{coh}}_8|$ below $1.5\times10^{-4}$, and the rate rises by $2.6\times10^{-4}$, which leaves $5.26\times10^3$ unchanged at three digits. The remaining uncertainties enter the power form \cref{eq:8.8-M} through $\ln(1 + \delta_k)$:
\begin{equation*}
  \begin{aligned}
  &\text{phonon frequencies, } +\SI{2.6}{\percent}: & \ln(1 + \delta) &= +2.4\times10^{-3}, \\
  &\text{anharmonic kinetic energy, } -\SI{0.14}{\percent}\ (\lambda_3 = 0.01, & &-8\times10^{-5}, \\
  &\qquad\lambda_4 = 2.5\times10^{-4}): & & \\
  &\text{anharmonic third cumulant (est.)}: & &+5\times10^{-5}.
  \end{aligned}
\end{equation*}
The kinetic-energy template is a dilation of each spectrum about $E_R$ by $(1 + \delta E^{\mathrm{kin}}/E^{\mathrm{kin}})^{1/2}$, which keeps the area and the mean and changes the second moment by the computed amount; the rate changes by 0.12 times the relative change of the width, and the extreme $-\SI{0.2}{\percent}$ of the range of \cref{sec:11.5} would give $-1.2\times10^{-4}$. For the third cumulant, the whole difference between the incoherent spectrum and a Gaussian of the same mean and variance changes the rate by \SI{1.9}{\percent}, and the anharmonic change of the cumulants is of relative order $\eta_{\mathrm{anh}}^2$. The anharmonic few-phonon weights barely enter: below \SI{20}{keV}, where they are large, $E_R$ is under \SI{8}{meV}, and those momentum transfers carry $3\times10^{-4}$ of the rate. With independent standard-normal $\vartheta_k$ the widths add in quadrature, \cref{eq:C.14-quadrature},
\begin{equation*}
  \sigma_{\ln\mathcal{R}} = \bigl[(2.4\times10^{-3})^2 + (8\times10^{-5})^2 + (5\times10^{-5})^2\bigr]^{1/2} = 2.4\times10^{-3} ,
\end{equation*}
and the rate lies between $5.25\times10^3$ and $5.27\times10^3$ per kilogram-year at one standard deviation, a factor of 1.0024 either way. The phonon frequencies carry all of it.

The free recoil gives $4.91\times10^3$, \SI{6.8}{\percent} lower, 29 band widths away. With the threshold raised to \SI{400}{meV}, near the endpoint \SI{0.48}{eV}, the crystal rate falls to \SI{110}{kg^{-1}.yr^{-1}}, half of it from momentum transfers below \SI{118}{keV}, $x_R = 6.5$. The free recoil gives 6.8, a sixteenth of it, because above \SI{400}{meV} most events come from the Doppler tail of the recoil spectrum, which the free nucleus does not have. The rate now samples the width of that spectrum, the dilation coefficient grows from 0.12 to 2.2, and the frequency template becomes $+\SI{3.1}{\percent}$; the band is a factor of 1.031. The anharmonic kinetic energy gives $-1.5\times10^{-3}$, still a twentieth of it.

Without perturbation theory the same rate would be far less certain. At \SI{80.4}{keV}, the median momentum transfer of the rate above \SI{100}{meV}, the five exact moments of the response allow any weight above threshold between 0.19 and 0.95 by the Markov--Krein band, against the harmonic 0.58.
\end{workedexample}

\section*{Summary}

A nucleon-coupled particle too light to give a nucleus a free recoil above threshold excites phonons, because the whole crystal takes up the momentum (\cref{sec:11.1}). At momentum transfers well inside the zone the one-phonon response per unit mass depends on the material only through the sound speeds and polarizations of the acoustic branches. The acoustic edge, set by the fastest longitudinal sound speed, \SI{6.1}{keV} for silicon at a \SI{1}{meV} threshold, lies 21 times above the free-target threshold mass. For a heavy mediator the rate reaches a plateau proportional to the polarization-weighted inverse sound speed, 654 events per kilogram and year in silicon at $\sigma_p = 10^{-40}\,\mathrm{cm^2}$.

A dark photon couples to charge and, in a polar crystal, excites the longitudinal optical phonon through the screened Born charge (\cref{sec:11.2}); for GaAs $Z^*/\varepsilon_\infty = 0.194$ from the measured splitting. Above the zone a nucleus leaves behind several phonons, Poisson distributed in the self part of the response, and the exact $\omega$-integrated weight, a static sum over pairs, corrects the incoherent approximation at every momentum transfer (\cref{sec:11.3}). At large momentum transfer the phonon orders merge into one peak at $E_R$ (\cref{sec:11.4}). At \SI{100}{keV} in silicon it lies at \SI{191}{meV} and is \SI{88}{meV} wide, and the free nuclear recoil with the Helm form factor, its limit, holds for a heavy particle within \SI{10}{\percent} from $x_R = 2.3$. In a crystal the free-nucleus formula for fermionic absorption on nuclei overstates the rate 16 times at \SI{10}{keV}.

Anharmonicity enters the rate first at second order, by time reversal (\cref{sec:11.5}); in silicon it changes the impulse width by less than a part in a thousand, and strongly anharmonic targets at detector temperatures have only the Markov--Krein band. Couplings that change from site to site give a broadband response that follows the density of states (\cref{sec:11.6}), and through the isotopes a baryon-minus-lepton vector is absorbed in silicon at least $5.0\times10^3$ times faster than through the electromagnetic channel at \SI{50}{meV}. A silicon crystal at $45^\circ$N modulates by \SI{0.9}{\percent} in its electronic rate and \SI{0.58}{\percent} in its lattice rate at \SI{2}{MeV} over the sidereal day (\cref{sec:11.7}). At neV to \si{\micro eV} masses the field is coherent over the body, and the acoustic modes of a resonator absorb it (\cref{sec:11.8}).

At \SI{30}{MeV} with a heavy mediator the silicon rate above \SI{100}{meV} is $5.26\times10^3$ events per kilogram-year, uncertain by a quarter of a percent, almost all of it from the computed phonon frequencies (\cref{sec:11.9}). The polar rates are set by $Z^*/\varepsilon_\infty$, which a computed $\varepsilon_\infty$ misses by enough to move a rate by \SI{37}{\percent}. Throughout the chapter the electrons stayed in their ground state and followed the nuclei. When a nucleus is kicked they do not follow at once, and \cref{chap:12} computes how often a nuclear recoil also excites them.
